\section{The value of the second flow}
\label{sec:value}

\subsection{Representation gain}

\begin{theorem}[Exact loss from suppressing one flow]
\label{thm:gap}
Under Assumption~\ref{ass:mean}, define
\begin{equation}
G(x)=\inf_{g\in\cF_1(x)}\E[(m_0(x,Y)-g(Y))^2\mid X=x].
\label{eq:gap_definition}
\end{equation}
If $v(x)=0$, then $G(x)=0$. If $v(x)>0$, then
\begin{equation}
G(x)=v(x)\min\left\{\frac{R_0(x)^2}{A(x)},\frac{U_0(x)^2}{B(x)}\right\}.
\label{eq:gap_result}
\end{equation}
Therefore $G(x)>0$ if and only if $v(x)>0$, $U_0(x)>0$, and $R_0(x)>0$.
\end{theorem}

The theorem gives a complete population answer for the comparator in Equation~\eqref{eq:f1}. One flow is exact when the conditional state is degenerate or one conditional-mean component is absent. Otherwise the two-flow class has a strict representation advantage. Let
\begin{equation}
\kappa(x)=\min\left\{\frac{R_0(x)^2}{A(x)},\frac{U_0(x)^2}{B(x)}\right\};
\label{eq:kappa}
\end{equation}
then $G(x)=v(x)\kappa(x)$. State dispersion supplies distinguishability, while $\kappa(x)$ measures coactivity: the cost of forcing one active direction to zero.

\subsection{Estimation cost and a single selection index}

To expose the additional estimation cost, fix states $y_1,\ldots,y_n$ and observe $Z=m+\varepsilon$ with $\varepsilon\sim\mathcal N(0,\sigma_\varepsilon^2 I_n)$. Let $P_2$ project onto the span of $(1-y_i)_{i=1}^n$ and $(-y_i)_{i=1}^n$. An oracle selects the better of the two one-dimensional spans, with projector $P_1$. Let $G_n=n^{-1}\|(I-P_1)m\|_2^2$.

\begin{proposition}[One additional degree of freedom]
\label{prop:benchmark}
If the two-column design has rank two, then
\begin{align}
\E\left[n^{-1}\|P_2Z-m\|_2^2\right]&=\frac{2\sigma_\varepsilon^2}{n},\nonumber\\
\E\left[n^{-1}\|P_1Z-m\|_2^2\right]&=G_n+\frac{\sigma_\varepsilon^2}{n},
\label{eq:benchmark_risks}
\end{align}
and
\begin{equation}
n^{-1}\{\|P_1Z-m\|_2^2-\|P_2Z-m\|_2^2\}
\overset{d}{=}G_n-\frac{\sigma_\varepsilon^2}{n}\chi_1^2.
\label{eq:benchmark_distribution}
\end{equation}
\end{proposition}

In the interior case, $G_n=\widehat v\min\{R_0^2/\widehat A,U_0^2/\widehat B\}$. The preceding theorem and proposition therefore meet in one dimensionless index:
\begin{equation}
\boxed{\begin{aligned}
\Gamma_n&:=\frac{nG_n}{\sigma_\varepsilon^2}=I_{\tau,n}\,\kappa_n,\\
I_{\tau,n}&:=\frac{n\widehat v}{\sigma_\varepsilon^2},\qquad
\kappa_n:=\min\left\{\frac{R_0^2}{\widehat A},\frac{U_0^2}{\widehat B}\right\}.
\end{aligned}}
\label{eq:gamma}
\end{equation}
Here $I_{\tau,n}$ is information in the turnover direction and $\kappa_n$ is coactivity. In the oracle-selected, span-relaxed Gaussian benchmark, the two-flow fit has lower expected fitted-mean error exactly when $\Gamma_n>1$, and
\begin{equation}
\Pp(\operatorname{Err}_2<\operatorname{Err}_1)=F_{\chi_1^2}(\Gamma_n).
\label{eq:win_probability}
\end{equation}
This is the classical one-extra-degree-of-freedom bias--variance calculation specialized to the opposing-flow basis \citep{mallows1973}. It is a calibrated local benchmark, not an exact threshold for bounded neural networks, heteroskedastic count fractions, or a dynamic panel. Its purpose is to identify the two quantities that any practical selection rule must balance.

\subsection{An exactly solvable case}
\label{sec:toy}

The supplied solvable case isolates the dispersion axis without changing the rates or noise level. At one driver value, let
\begin{align}
Y&\in\left\{\frac12-\delta,\frac12+\delta\right\},\nonumber\\
\Pp\left(Y=\frac12-\delta\right)
&=\Pp\left(Y=\frac12+\delta\right)=\frac12.
\label{eq:two_state_design}
\end{align}
with $U\ge R>0$. Then $A=B=1/4+\delta^2$, $C=1/4-\delta^2$, $v=\delta^2$, and
\begin{equation}
G(\delta)=\frac{4R^2\delta^2}{1+4\delta^2},
\qquad
\Gamma(\delta)=\frac{4nR^2\delta^2}{\sigma_\varepsilon^2(1+4\delta^2)}.
\label{eq:toy_closed_form}
\end{equation}
For an even balanced sample,
\begin{equation}
\E[\operatorname{Err}_1-\operatorname{Err}_2]
=\frac{4R^2\delta^2}{1+4\delta^2}-\frac{\sigma_\varepsilon^2}{n},
\label{eq:toy_risk_difference}
\end{equation}
so the expected-risk crossover is
\begin{equation}
\delta_c=\frac{\sigma_\varepsilon}{2\sqrt{nR^2-\sigma_\varepsilon^2}},
\label{eq:delta_critical}
\end{equation}
whenever the right-hand side lies in $(0,1/2)$. For $U=0.20$, $R=0.12$, $n=40$, and $\sigma_\varepsilon=0.15$, Equation~\eqref{eq:delta_critical} gives $\delta_c=0.1008$. Fifty thousand Gaussian replications at each grid point reproduce the closed-form risks. The case deliberately isolates dispersion and attains the upper endpoint $\lambda_{\min}(Q)=2v$; the real outage design in Section~\ref{sec:real_geometry} lies close to the lower endpoint $\lambda_{\min}(Q)=v$. Panel~(c) of Figure~\ref{fig:toy} adds the complementary coactivity slice: at fixed dispersion, the gap vanishes when either component is absent.

\begin{figure*}[t]
\centering
\includegraphics[width=0.98\textwidth]{figures/flow_selection_solvable_case.pdf}
\caption{Exactly solvable flow-selection case. The first two panels reproduce the supplied dispersion experiment and its finite-sample crossover; the third isolates the coactivity factor. In panels (a)--(b), the vertical line marks $\Gamma=1$, equivalently $\delta=0.1008$ under the stated parameters; panel (c) marks the one-flow branch switch at $R/U=1$.}
\label{fig:toy}
\end{figure*}

\begin{table}[t]
\centering
\caption{Selected points from the exactly solvable case. Risks are expected fitted-mean MSE.}
\label{tab:toy}
\small
\begin{tabular}{@{}ccccc@{}}
\toprule
$\delta$ & $\Gamma$ & One flow & Two flows & Lower risk\\
\midrule
0.0494 & 0.247 & 0.000702 & 0.001125 & One\\
0.1008 & 1.000 & 0.001125 & 0.001125 & Tie\\
0.1963 & 3.417 & 0.002485 & 0.001125 & Two\\
\bottomrule
\end{tabular}
\end{table}

\subsection{The same criterion under event shift}

The best one-flow projection depends on the conditional state distribution. Let $P$ and $Q$ denote source and target environments, and suppose locally that $U_0(x)$ and $R_0(x)$ are invariant while the law of $Y\mid X=x$ changes. For the interruption-only branch, the source population coefficient is $a_P^*=U_0-R_0C_P/A_P$. If the source and target optima are interior, its target approximation risk is
\begin{equation}
\cR_Q^U(a_P^*)=R_0^2\frac{v_Q}{A_Q}
+R_0^2A_Q\left(\frac{C_P}{A_P}-\frac{C_Q}{A_Q}\right)^2.
\label{eq:projection_shift}
\end{equation}
The second term is a projection-shift penalty. Under the same local homoskedastic benchmark, estimating two flows from $n_P$ source observations incurs expected target excess risk
\begin{equation}
\frac{\sigma_\varepsilon^2}{n_P}
\left[1+\frac{v_Q+(\mu_Q-\mu_P)^2}{v_P}\right].
\label{eq:two_transfer_cost}
\end{equation}
Thus event transfer modifies both sides of the same comparison: one-flow approximation gains a shift penalty, whereas turnover estimation gains a leverage penalty. When $P=Q$, the comparison reduces to $\Gamma_n>1$. We do not use a plug-in transfer score as an empirical selector because its predictive validity has not yet been established.
